% Generated draft: requires mathematical and source review.
\documentclass[11pt]{article}
\usepackage[margin=25mm]{geometry}
\usepackage{fontspec}
\setmainfont{Latin Modern Roman}
\usepackage{amsmath,amssymb,mathtools,unicode-math}
\setmathfont{Latin Modern Math}
\usepackage{enumitem,needspace,xurl,hyperref}
\setlength{\parindent}{0pt}
\setlength{\parskip}{5pt}
\setlength{\emergencystretch}{3em}
\newcommand{\N}{\mathbb{N}}
\newcommand{\Z}{\mathbb{Z}}
\newcommand{\Q}{\mathbb{Q}}
\newcommand{\R}{\mathbb{R}}
\newcommand{\C}{\mathbb{C}}
\newcommand{\F}{\mathbb{F}}
\newcommand{\A}{\mathbb{A}}
\newcommand{\PP}{\mathbb{P}}
\newcommand{\E}{\mathbb{E}}
\newcommand{\eps}{\varepsilon}
\newcommand{\dd}{\,\mathrm{d}}
\begin{document}
\section*{UnsolvedMath ID: 30004091}
\textbf{Source problem code:} catalogue entry
\subsection{Definitions and mathematical statement}
Put \(K=\C((t))\) and \(R=\C[[t]]\). A special fiber requires a degeneration model over \(R\), rather than a variety over \(K\) alone. Accordingly, interpret \(\mathcal X\) as such a model, with generic fiber \(X_\eta\) over \(K\), special fiber \(X_0\), and reduced special fiber \(D=(X_0)_{\mathrm{red}}\). This interpretation of the base is implicit in the source and requires confirmation.

In the usual irreducible convention, a compact hyperkähler manifold is a simply connected compact Kähler manifold whose space of holomorphic two-forms is generated by an everywhere nondegenerate form. The source does not specify whether this irreducible convention is required, or how its complex-geometric terminology is implemented over \(K\).

The abbreviation dlt means divisorial log terminal. For a dlt pair \((\mathcal X,D)\) with reduced boundary, its dual complex \(\mathcal D(D)\) has a vertex for each irreducible component of \(D\), and cells corresponding to its intersection strata, with attaching relations determined by incidence. In particular, distinct strata can give distinct cells even when they involve the same components; one must not replace this complex by the abstract nerve of the components. Write \(\lvert\mathcal D(D)\rvert\) for its underlying topological space. The source does not explicitly specify the boundary or special-fiber multiplicity convention.

A good minimal model ordinarily has a relatively nef and semiample log canonical divisor for the relevant pair. The source retains the hypothesis \emph{good minimal dlt degeneration} without specifying that pair or its precise relative canonical-divisor conditions. Likewise, \emph{maximally unipotent} is retained as a hypothesis: the record does not identify the cohomological degree, monodromy operator, or nilpotency condition used to define it.

\textbf{Question, with unresolved conventions made explicit.} For every maximally unipotent good minimal dlt degeneration \(\mathcal X\) of compact hyperkähler manifolds over \(\C((t))\), is there a homeomorphism
\[
\lvert\mathcal D(D)\rvert\ \cong\ \PP^n(\C)?
\]
Here \(\PP^n(\C)\) is the space of complex lines in \(\C^{n+1}\), with its usual topology. The source leaves \(n\) undefined. The expected dimensional interpretation is \(\dim X_\eta=2n\), in which case the proposed target has real dimension \(2n\); this interpretation is provisional and must be confirmed. The question therefore remains conditional on resolving the dimensional and degeneration conventions above.
\subsection{Short English statement}
The question asks whether every maximally unipotent good minimal dlt degeneration of compact hyperkähler manifolds has special-fiber dual complex homeomorphic to complex projective space. The intended relation is presumably \(\dim X_\eta=2n\), but the source does not define \(n\) or fully specify the degeneration and monodromy conventions. This is a statement-only draft requiring human mathematical review.
\subsection{Sources}
\begin{itemize}
\item Local UnsolvedMath registry, permanent ID \(30004091\), title \emph{Dual Complexes of Maximally Degenerate Hyperkähler Varieties}, exact source problem code \texttt{catalogue entry}, in the Oberwolfach Reports Open Problems collection. No individual proposer is supplied; the proposed year is recorded as 2019.
\item Supplied source citation: \emph{Non-Archimedean Geometry and Applications} (2020), Oberwolfach Reports, pp.\ 513--575, DOI \texttt{10.4171/owr/2019/8}. Supplied report link: \url{https://doi.org/10.4171/owr/2019/8}. The citation year is 2020, whereas the proposed year and DOI contain 2019; both are retained without reconciliation.
\item Supplied literature reference: \emph{Degenerations of hyper-Kahler manifolds}, \url{https://aif.centre-mersenne.org/item/AIF_2018__68_7_2837_0/}.
\item Supplied literature reference: \emph{Oberwolfach report}, \url{https://doi.org/10.4171/owr/2019/8}.
\item Literature assessment dated exactly \texttt{2026-08-22}: the supplied curation labels the problem open and states that, for type-III hyperkähler degenerations, the dual complex is known to be a closed pseudomanifold with the rational homology of complex projective space, while the stronger homeomorphism claim remains open beyond special cases. This is supplied dated curation, not an independently verified literature assessment; the record explicitly requests independent verification of its time-sensitive status and cited claims.
\item The supplied statement status is \texttt{corrected\_verified}. Its verification note reports that the official Oberwolfach report was checked directly and that the clean statement preserves the target while restoring definitions and removing source-dependent wording. This is registry provenance, not a claim of source consultation or mathematical verification in this draft.
\end{itemize}
\end{document}
